Throughout $p = p_{k}$, $\Omega = \Omega_{k}$, $\tau = \tau_{\epsilon}$, and $p' = \prod_{j}v_{j}w_{j}$. The relations give $(w_{j}v_{j})^{2} = w_{j}(1 - w_{j}v_{j})v_{j} = w_{j}v_{j}$ and likewise $(v_{j}w_{j})^{2} = v_{j}w_{j}$; the elements $w_{j}v_{j}$ are even and lie in different planes, so they commute with each other and with $v_{i}, w_{i}$ for $i \ne j$; hence $p$ and $p'$ are idempotents, and $w_{j}p = 0$ because $w_{j}w_{j}v_{j} = Q(w_{j})v_{j} = 0$. Further $v_{j}(w_{j}v_{j}) = (1 - w_{j}v_{j})v_{j} = v_{j}$ and $(v_{j}w_{j})v_{j} = v_{j}$, hence
\[
\Omega p = \Omega = p'\Omega, \qquad \tau(p) = p'
\]
(the reverse of $w_{j}v_{j}$ is $v_{j}w_{j}$, the factors commute, and $p$ is even).

\emph{Relation to the stored pairing.} Proposition \claimref{std-algebra.fock-representation} gives the Fock space as the minimal left ideal on a vacuum idempotent with its occupation basis, and Proposition \claimref{std-algebra.chevalley-pairing} proves, for that model, a nondegenerate form pairing grade $m$ only with grade $k - m$, with symmetry sign $(-1)^{k(k-1)/2}$ and every vector self-adjoint; there the creators are the $w_{i}$ and the idempotent is $v_{1}w_{1}\cdots v_{k}w_{k}$, which is the present setting with the two null families exchanged. By the uniqueness proved in (i) below, $\beta_{+}$ is that pairing up to a factor. The proofs of (i) to (iii) are written out in the present notation because $\beta_{-}$, uniqueness and the signature are needed and are not stored.

\emph{The module.} $B_{k}$ is spanned by the products $v_{S}w_{T}$, and $w_{T}p = 0$ unless $T$ is empty, so $F_{k} = B_{k}p$ is spanned by the $v_{S}p$. The map $v_{S} \mapsto v_{S}p$ from the exterior algebra of $W_{k} = \operatorname{span}\{v_{j}\}$ is a nonzero map of $B_{k}$-modules for the Chevalley action ($v_{j}$ by multiplication, $w_{j}$ by the graded derivation with $w_{j}\cdot v_{i} = \delta_{ij}$, which is how $w_{j}$ moves through $v_{S}$ to reach $p$), and that module is irreducible of dimension $2^{k}$ with $B_{k}$ acting as its full endomorphism algebra (Theorem \claimref{tower.fock-module-one-qubit-per-rung}); so the map is an isomorphism and the $v_{S}p$ are a basis.

\emph{(i).} For $\psi = xp$, $\phi = yp$: $\tau(\psi)\phi = p'\,\tau(x)y\,p \in p'F_{k}$. Now $p'v_{S}p = 0$ unless $S$ is the full set: if $j \notin S$, the factor $v_{j}w_{j}$ of $p'$ can be moved next to $v_{S}p$ and $w_{j}v_{S}p = \pm v_{S}w_{j}p = 0$. For the full set $p'\Omega p = \Omega$. So $p'F_{k}$ is the line of $\Omega$, and $\beta_{\epsilon}$ is defined and bilinear. On the basis, $\tau(v_{S}p)v_{T}p = p'\,\tau(v_{S})v_{T}\,p$ and $\tau(v_{S})v_{T} = \pm v_{S}v_{T}$ reordered, which vanishes if $S$ and $T$ meet ($v_{j}^{2} = Q(v_{j}) = 0$), is killed by $p'$ if they do not cover all rungs, and is $\pm\Omega$ if $T = S^{c}$. So the Gram matrix has exactly one nonzero entry, $\pm 1$, in each row and column, and $\beta_{\epsilon}$ is nondegenerate. Since $\tau$ is an antiautomorphism with $\tau(a) = \epsilon a$ on vectors, $\tau(a\psi)\phi = \tau(\psi)\tau(a)\phi = \epsilon\,\tau(\psi)a\phi$, which is the adjoint property. If $\beta'$ is another form with the property, nondegeneracy of $\beta_{\epsilon}$ gives a linear map $L$ of $F_{k}$ with $\beta'(\psi, \phi) = \beta_{\epsilon}(\psi, L\phi)$, and for every vector $a$
\[
\beta_{\epsilon}(\psi, La\phi) = \beta'(\psi, a\phi) = \epsilon\,\beta'(a\psi, \phi) = \epsilon\,\beta_{\epsilon}(a\psi, L\phi) = \beta_{\epsilon}(\psi, aL\phi),
\]
so $L$ commutes with every vector, hence with $B_{k}$, which acts as the full endomorphism algebra of $F_{k}$; $L$ is a scalar.

For $\beta_{-}$: the grade involution $\alpha$ of $B_{k}$ preserves $F_{k}$ ($p$ is even) and acts on $v_{S}p$ by $(-1)^{|S|}$; since $\tau_{-} = \tau_{+}\alpha$, $\tau_{-}(\psi)\phi = \tau_{+}(\alpha\psi)\phi$, that is $\beta_{-}(\psi, \phi) = \beta_{+}(\alpha\psi, \phi)$.

\emph{(ii).} Apply $\tau$ to $\tau(\psi)\phi = \beta_{\epsilon}(\psi, \phi)\Omega$: the left side becomes $\tau(\phi)\psi = \beta_{\epsilon}(\phi, \psi)\Omega$ ($\tau$ is an involution), and $\tau(\Omega) = \epsilon^{k}v_{k}\cdots v_{1} = \epsilon^{k}(-1)^{k(k-1)/2}\Omega$, the reverse of a product of $k$ anticommuting vectors. Hence $\beta_{\epsilon}(\phi, \psi) = \epsilon^{k}(-1)^{k(k-1)/2}\beta_{\epsilon}(\psi, \phi)$. The exponent $k(k-1)/2$ is even exactly for $k \equiv 0, 1 \pmod 4$, and $k(k-1)/2 + k = k(k+1)/2$ exactly for $k \equiv 0, 3$, which gives the table and the two existence statements.

\emph{(iii).} Let $\beta_{\epsilon}$ be symmetric. Since $k \ge 1$, $S \ne S^{c}$, so every basis state is null, and by (i) the plane of $v_{S}p$ and $v_{S^{c}}p$ is hyperbolic and orthogonal to every other basis state. There are $2^{k-1}$ such planes, so the signature is $(2^{k-1}, 2^{k-1})$ and the Witt index $2^{k-1}$; Theorem \claimref{tower.witt-index-quantum-capacity} identifies the Witt index with the largest number of mode pairs realised by vectors. For any nondegenerate quadratic form on a space of dimension $2^{k}$ the Witt index is at most half the dimension, which is the last statement. If $\beta_{\epsilon}$ is antisymmetric the same planes are symplectic planes. The form pairs a state of grade $|S|$ only with one of grade $k - |S|$; for odd $k$ these have opposite parity, so each parity half is totally null and the form pairs the halves, and for even $k$ they have equal parity, so the halves are orthogonal and each is nondegenerate.

\emph{(iv).} $s_{+}(3) = (-1)^{3} = -1$ and $s_{-}(3) = (-1)^{3}(-1)^{3} = +1$; (iii) gives signature $(4,4)$ and the null halves, each of dimension $2^{2} = 4$. The vector space of $B_{4} = \Cl(4,4)$ is the orthogonal sum of the four hyperbolic planes of its null pairs, and $(F_{3}, \beta_{-})$ is the orthogonal sum of four hyperbolic planes by (iii), so the two are isometric, as any two hyperbolic spaces of equal dimension are. No isometry intertwines the actions of $B_{3}$. The bivectors of $B_{3}$ act on $F_{3}$ by left multiplication; under them $F_{3}$ is the sum of its two parity halves, and a vector $\psi$ annihilated by all of them is zero: $v_{i}w_{j}\psi = 0$ and $w_{i}v_{j}\psi = 0$ for $i \ne j$ and $(v_{j}w_{j} - w_{j}v_{j})\psi = 0$ force every coefficient of $\psi$ on the basis $v_{S}p$ to vanish (the last relation multiplies $v_{S}p$ by $+1$ if $j \in S$ and $-1$ if not, never by zero; the check computes the common kernel as zero). On the vector space of $B_{4}$, spanned by the vectors of $B_{3}$ and $v_{4}, w_{4}$, the bivectors of $B_{3}$ act by commutator and commute with $v_{4}$ and $w_{4}$, so they annihilate that plane. An intertwining isometry would carry the plane to nonzero annihilated vectors of $F_{3}$, which do not exist. The rule $k \mapsto 2^{k-1}$ gives $8$ at $k = 4$, where the tower's next floor has $5$ pairs. For the equality: $2^{k-1} - (k+1)$ is $-1, -1, 0, 3$ at $k = 1, 2, 3, 4$ and increases by $2^{k-1} - 1 > 0$ at each step from $k = 2$, so it vanishes only at $k = 3$. $s_{+}(2) = -1$ and $s_{-}(2) = (+1)(-1) = -1$, and by the uniqueness in (i) there is no other form for which the vectors are self-adjoint or anti-self-adjoint.

\emph{(v).} The first rule and its counts are Theorem \claimref{synthesis.exponential-reading-of-the-tower-is-the-iterated-fock-functor-not-the-ascent}(i). The second rule is (iii): a symmetric $\beta_{\epsilon}$ makes $F_{k}$ a hyperbolic space of Witt index $2^{k-1}$, whose Clifford algebra is $\Cl(2^{k-1}, 2^{k-1}) = B_{2^{k-1}}$; $k \mapsto 2^{k-1}$ gives $4$, $8$, $128$, $2^{127}$ from three. At $k = 3$ it must use $\beta_{-}$, the only symmetric form there by (ii) and the uniqueness in (i); every term from the second on is a power of two at least four, so it is divisible by four and both forms are symmetric there by (ii), with the same count; at $k = 2$ no symmetric form exists by (iv). The state counts are $2^{k}$ at $k = 3, 4, 8, 128$, and $3 \times 10^{38} < 2^{128} < 4 \times 10^{38}$. For the horizon count, the bracket $2^{407} < |J| < 2^{408}$ and the rider that sizes only are compared are quoted from Theorem \claimref{synthesis.exponential-reading-of-the-tower-is-the-iterated-fock-functor-not-the-ascent}(v); given it, compare base-two logarithms: $127 < 407$ and $408 < 2^{127} - 1$. The proved growth is one pair per rung (Theorem \claimref{tower.witt-index-quantum-capacity}, and Theorem \claimref{synthesis.exponential-reading-of-the-tower-is-the-iterated-fock-functor-not-the-ascent}(ii) for the ascent from three): $k \mapsto k + 1$ is neither rule, and its value at three equals the second rule's as a number only, by (iv). That the second rule is a reading of the stored phrase is a judgement. The sentence that no stored result selects a rule is argued: the tower results in the store were searched for a statement relating the module of one rung to the generating vectors of a later one, and none was found; this is not an audit of a full cone.
